\section{The transverse frame estimate}\label{f:section}

Write
\[
 q=\frac{2(k+2)}k,\qquad p_k=\frac{2(k+1)}k,\qquad m=k+1.
\]
All packet coordinates, envelopes, and weighted packet sizes in this
section are those of Section~\ref{p:section}. In particular, an inverse
frequency width $r$ corresponds to spatial width $r$, time length $r^2$,
and spacetime volume $r^{k+2}$. A packet mask means multiplication by a
selected sum of the smooth packet cutoffs at its own scale. Such sums
have bounded absolute value. A restriction of a mask is allowed to keep
any subset of its terms.

\Needspace{36\baselineskip}
\begin{theorem}[Frame estimate]\label{f:frame}
Fix $k\ge2$, $c>0$, and fixed packet bounds $C_0$. There is a finite
integer $P_0$ such that, for every fixed $P\ge P_0$, there are a finite
integer $J_0$ and a constant $C$, depending only on $k,c,C_0,P$ and the
packet conventions, with the following property. Let $h\ge2$, let
$V_1,\ldots,V_m$ be scale-$h$
packet sums on one time interval of length at most $C_0h^2$, and suppose
their normalized compact profiles have derivatives through order $J_0$
bounded by $C_0$. Let $\Xi_i$ be their assigned frequency-label sets,
with the fixed cap enlargements included. Assume bounded phase-space
counting and
\[
 \big|\det((2\xi_1,1),\ldots,(2\xi_m,1))\big|\ge\delta>0
 \quad(\xi_i\in\Xi_i).
\]
Use $P$ for the packet envelope exponent. Let $Y$ be a set of
distinct unit cubes and, for every $Q\in Y$ and $i$, let $z_{i,Q}\in Q$
satisfy
\[
 |V_i(z_{i,Q})|\ge A_i\ge h^c,
 \qquad B_i(z_{i,Q})^2=\sum_Ts_{i,T}^2w_T(z_{i,Q})\le1.
\]
The points $z_{i,Q}$ need not agree for different $i$. Set
\[
 H_i=A_i^{2/k},\qquad H_g=\Big(\prod_{i=1}^mH_i\Big)^{1/m},
 \qquad D_i=h^{k+2}\sum_Ts_{i,T}^2.
\]
Then
\begin{equation}\label{f:mainbound}
 |Y|\le C\delta^{-C}H_g^{-2}\sum_{i=1}^m\frac{D_i}{A_i^2}.
\end{equation}
Only the finite number $J_0$ of profile derivatives is required.
\end{theorem}

Unit translations change the square envelopes by a fixed factor.
Consequently we use the cube centers for counting and retain the
individual witnesses when testing the fields. Every local norm estimate
below uses a weighted supremum on the relevant cube; thus none of the
arguments replaces separate witnesses by an unsupported common
witness.

The proof is by contradiction. A failing configuration determines
transition scales and an excess packet amplitude. We encode the
amplitude ratios along a finite hierarchy by logarithmic profiles.
Projection estimates force each profile to have only a zero slope or
the full Cauchy slope, and an incidence argument locates the interval
on which the full slope occurs. The resulting transition ratio is
subpower, contradicting the original excess amplitude and its energy
budget. Smoothing is needed only when that excess grows
subpolynomially in the original packet scale.

\subsection{Exact packets and a normalized countersequence}

\begin{lemma}[Reduction to exact packets]\label{f:exact}
It suffices to prove Theorem~\ref{f:frame} for translated exact
paraboloid extensions, with smooth data supported in bounded native
frequency boxes, multiplied by a common Schwartz function of $t/h^2$
which is positive on the test interval and has compact Fourier support.
All required seminorms of these exact packets can be fixed in advance.
\end{lemma}

\begin{proof}
Use a fixed enlargement of the test interval and a cutoff equal to one
there. Fourier transform a compact native profile in its spatial
variables and insert a smooth partition into unit frequency boxes
centered at $a\in\Z^k$. In the $a$th box remove the quadratic evolution
factor, and expand the remaining smooth function of native time in a
Fourier series on the enlarged interval. This writes the profile as a
sum, indexed by $(a,n)\in\Z^{k+1}$, of exact extensions times common
time modulations. Integration by parts in the Fourier coefficient
formulas shows that, for prescribed finite $J,M$,
\[
 \|\text{data}_{a,n}\|_{C^J}
 \le C_{J,M}(1+|a|+|n|)^{-M},
\]
provided finitely many original derivatives are bounded. The physical
frequency center changes by $a/h$. All intercepts are specified at one
common reference time. Recentring about the shifted trajectory on the
test interval changes a polynomial envelope by at most a fixed
polynomial in $1+|a|$.

For completeness, normalize each layer by
$c_{a,n}=(1+|a|+|n|)^{-M/2}$ and allocate threshold fractions proportional
to $(1+|a|+|n|)^{-M/4}$. Their sum is finite when $M$ is large. If a sum
has magnitude at least $A_i$, one of its normalized layers exceeds the
corresponding allocated threshold divided by $c_{a,n}$. Apply the exact
packet estimate to the tuples of these layers. In its $i$th energy term
the threshold dependence is
\[
 D_i A_i^{-2}\prod_{j=1}^m A_j^{-4/(km)}.
\]
Every threshold has a strictly negative power; hence, after taking $M$
large, the resulting mode sums converge, including the polynomial
envelope losses. Modes $|a|>h^{1/2}$ contribute less than any specified
negative power of $h$ by weighted Cauchy and the same decay. For the
remaining modes transversality is preserved when
$\delta\ge h^{-\eps_0}$ and $\eps_0>0$ is sufficiently small. If
$\delta<h^{-\eps_0}$, the polynomial bound in $h$ from the elementary
packet estimate is absorbed by $\delta^{-C}$. This also disposes of
the discarded tail.

An exact packet outside its central time bin has only polynomial
growth in native time in each of the finitely many spatial localization
seminorms: differentiate the compactly supported frequency integral and
integrate by parts. Multiplication by a common positive-on-the-slab
Schwartz cutoff makes all relevant spacetime integrals finite. Choose
its decay after those seminorms. The whole argument therefore uses a
fixed finite derivative order.
The initial native-time cutoff used for the Fourier expansion is
distinct from the final compact-Fourier multiplier
$\chi_h(t)=\chi((t-t_0)/h^2)$. The latter is bounded above and below
by positive fixed constants on the test slab. Multiplying the
exact-layer sum by it therefore changes the witness thresholds and
packet bounds by fixed factors, which suffices for the reduction.
\end{proof}

Suppose now that the exact-packet assertion fails. By choosing the
successive counterexamples to violate every fixed constant and every
fixed inverse power of $\delta$, we obtain $h\to\infty$ with
$\delta=h^{-o(1)}$. Absolute Cauchy gives $H_i\lesssim h$. Proposition
\ref{p:weighted} gives
\[
 |Y|\le h^{o(1)}D_iH_i^{-k-2}.
\]
Comparing this for all $i$ with the failed estimate proves
\begin{equation}\label{f:hbalance}
 H_i/H_g=h^{o(1)},\qquad |Y|H_i^{k+2}/D_i=h^{o(1)}.
\end{equation}
Here and below a positive quantity equal to a base to the power $o(1)$
is bounded both above and below in that sense.

Write the original exact-packet field as $V_i=\chi_h U_i$, where
$U_i$ is its uncut exact extension, with any common temporal
modulation removed. For the smooth inverse-width-$H_i$ cap projection
$P_b$, put $U_b=P_bU_i$ and $V_b=\chi_h U_b$. Thus $U_b$ always
denotes the full original cap extension before any transition masks.
Localize each actual cap field $V_b$ at inverse width $H_i$, and select
one dyadic packet amplitude class in each field:
\begin{equation}\label{f:original-amplitudes}
 a_T\asymp S_iH_i^{-k/2},\qquad S_i\gtrsim1.
\end{equation}
These are the \emph{original transition masks}. They will remain in
every later hierarchy, even though the fields inside them change.
The selected sum has high threshold
$A_i(1+|\log S_i|)^{-C}$ on a subset, again denoted $Y$, satisfying
\begin{equation}\label{f:selected-counter}
 |Y|\ge (\log(2+S))^{-C}H_g^{-2}\sum_iD_iH_i^{-k},
 \qquad S=\max_iS_i,
\end{equation}
with an arbitrarily large additional constant and inverse power of
$\delta$ along the countersequence. This countable pigeonhole is
legitimate: classes below a fixed small value of $S_i$ have a summably
small absolute contribution, and the threshold fractions above are
summable in the dyadic index. At a transition cell, local Cauchy gives
$a_T\lesssim(h/H_i)^{k/2}b_T$, where $b_T$ is its full descendant square
weight. The weighted square sum at witnesses is bounded. Consequently
classes larger than a fixed polynomial in $h$ have a negligible root
contribution, so that $\log S=O(\log h)$ for the selected tuple.

Spatial almost orthogonality, followed by the local weighted Bernstein
bound, gives
\begin{equation}\label{f:original-energy}
 H_i^{k+2}\sum_{T\text{ selected}}a_T^2\lesssim D_i.
\end{equation}
Bounded counting and \eqref{f:original-amplitudes} also give square
background $O(S_i)$ for the selected packets without using the original
$B_i$. \Cref{p:near-cauchy}, applied to these packets, yields
\begin{equation}\label{f:Sbalance}
 H_i/H_g+H_g/H_i+D_i/(|Y|H_i^{k+2})
       +|Y|H_i^{k+2}/D_i\lesssim S^C.
\end{equation}
It follows that $S\to\infty$ and $\delta=S^{-o(1)}$. Indeed, bounded
$S$, or a fixed positive value of $\log_S\delta^{-1}$, would make this
estimate contradict the choice of the countersequence in
\eqref{f:selected-counter}.

Pass to a subsequence in one of two cases. In the \emph{power case},
$\log S\gtrsim\log h$ and the logarithmic base is $\Lambda=h$. In the
\emph{subpower case}, $S=h^{o(1)}$ and the base is $\Lambda=S$. The next
subsection is needed only in the latter case.

The transition cells chosen in \eqref{f:original-amplitudes} will stay
fixed, even when we replace the fields localized by them.  We retain
their original amplitudes as data, writing
\[
 a_T^{\mathrm{orig}}\asymp S_iH_i^{-k/2}.
\]
Later assemblies have their own actual amplitudes; these need not remain
comparable to $a_T^{\mathrm{orig}}$.  The original energy bound
\eqref{f:original-energy}, however, continues to charge every surviving
transition cell.

Here is the count we seek.  We will show that the high set requires at
least
\[
 |Y|H_i^k\Lambda^{-o(1)}
\]
distinct surviving transition cells in each field.  Applied to a field
with $S_i=S$, their original amplitudes would give
\[
 D_i\gtrsim |Y|H_i^{k+2}S^2\Lambda^{-o(1)}.
\]
The scale balances will exclude this inequality in both cases.  To
obtain the count, we first build finite scale hierarchies through the
fixed transition masks.  Their selected amplitude ratios define
logarithmic profiles.  The projection and incidence arguments will
force the ratio at the transition scale to be subpower and bound the
number of visits to each cell.  Cauchy--Schwarz gives many incident cells at each high point, and
double counting yields the displayed global bound.  In the subpower case,
we first need a square background controlled in the base $S$; this is
the purpose of the smoothing argument that follows.

\subsection{Smoothing and clipping on the subpower scale}

Fix a field and abbreviate $H=H_i$, $D=D_i$. First remove original
transition packets visited by more than $S^{C_1}$ witnesses in a fixed
large power enlargement. Their squared-amplitude sum is at most
$|Y|S^{-C_1+C_2}$: count visits, using the fixed amplitude class and
bounded counting, and put the enlargement costs into $C_2$.
\Cref{p:near-cauchy}, with square background $O(S)$, shows that
their root contribution is negligible when $C_1$ is sufficiently large.
Each remaining high point needs at least $H^k/S^C$ nearby cap slots,
by absolute summation and \eqref{f:original-amplitudes}. Packet tails
outside the chosen enlargement are negligible by the fixed rapid
cutoff decay. Counting these incidences in a ball gives
\begin{equation}\label{f:ballcount}
 \#(Y\cap B_r)\lesssim S^C(r/H)^k(1+r/H^2),\qquad r\ge H.
\end{equation}
For each of the $O(H^k)$ cap directions, at most
$O((r/H)^k(1+r/H^2))$ packet cells meet the ball, up to the fixed power
enlargement. Their occupancy is at most $S^{C_1}$. Division by the
$H^k/S^C$ incidences required at each point proves the formula.

\begin{lemma}[Clipped smoothing]\label{f:smoothing}
Fix $C_3$ sufficiently large. There is a fixed $L$ such that, with
$K_0=HS^L\le h$, each cap field $V_b$ admits a scale-$K_0$ packet
representation with integrated square energy at most $CD$ and
envelopes
$e_b^2=\sum_{\gamma\in b}s_\gamma^2w_\gamma^{\rm work}$ satisfying
\begin{equation}\label{f:clipped}
 \sum_b\min(e_b^2(z),M_c)\le S^{o(1)},
 \qquad M_c=H^{-k}S^{C_3},
\end{equation}
outside $o(|Y|)$ witnesses. The work weight
$w_\gamma^{\rm work}=(1+d_\gamma)^{-P_{\rm work}}$ has a fixed order
$P_{\rm work}>P$, chosen in the proof; all subsequent levels of this
subpower tree use that order. The envelope and profile orders are
finite and independent of any subsequent scale-grid length. If
$K_0>h$, the tree instead ends at the unchanged original scale-$h$
packets with their original order $P$; their square envelope is
already bounded at the witnesses.
\end{lemma}

\begin{proof}
\textbf{The positive smoothing functional.}
We give separately the positive smoothing estimate and its comparison
with packet sizes. The original selections and deletions are left
unchanged and applied externally, and $\chi_h$ remains in every actual
field. Only the auxiliary positive density below uses the uncut
original cap field $U_b$. Fix a sufficiently large integer $J\ge1$,
independent of all later grids, and put
\[
 \rho=H^{k+2}\sum_{z\in Y}\delta_z,\qquad
 F_b(x,t)=w_h(t)|U_b(x,t)|^2,\qquad
 w_h(t)=\left(1+\frac{(t-t_0)^2}{h^4}\right)^{-J}.
\]
Spatial unitarity and the original fine-packet Gram bound give
\begin{equation}\label{f:auxiliary-energy}
 \sum_b\int F_b
 =c_Jh^2\sum_b\|U_b(t_0)\|_2^2
 \lesssim h^{k+2}\sum_Ts_T^2=D.
\end{equation}
Here $c_J=\int_{\R}(1+s^2)^{-J}\,\dd s<\infty$.
Indeed the cap projections have bounded overlap, and integration by
parts in the original normalized frequency data gives summable Gram
rows for the fine packets, using finitely many prescribed localization
seminorms. This charges the original data before every transition
mask. No division by $\chi_h$ is used. Choose an even
$Q$ sufficiently large compared with the original envelope order $P$ and
$k$, and put $j_K(x)=c_QK^{-k}(1+|x/K|^2)^{-Q/2}$.

Convolve in space with heat variance $K^2$ and with independent
symmetric jumps having density $j_r$ and intensity $\dd r/r$ at scale
$r\le K$. The limiting small-jump distribution exists because
$\int_0^K r^2\dd r/r<\infty$. Denote the resulting kernel by $J_K$.
It has
\begin{equation}\label{f:kernel}
 C^{-1}j_K\le J_K\le Cj_K,\quad
 |\nabla J_K|\le CK^{-1}j_K,\quad
 |\widehat J_K(\xi)|\le e^{-cK^2|\xi|^2}.
\end{equation}
Here is a useful direct verification of the polynomial tails. For the
lower bound request one jump in $[K/2,K]$ and all other displacement
at most $CK$, an event of fixed positive probability by the second
moment bound. The density of that jump is comparable to $j_K$ after
such a displacement. For the upper bound group jumps into
$[2^{-j-1}K,2^{-j}K]$, with Poisson counts $m_j$. If their total plus
the Gaussian has magnitude comparable to $|x|\gg K$, either the
Gaussian has such magnitude or one of the jumps in group $j$ has
magnitude at least $c2^{-j/2}|x|/(1+m_j)$. Conditioning on the other
variables and summing candidates bounds its density by
\[
 C(1+m_j)^{Q+1}K^{-k}2^{jk-jQ/2}(K/|x|)^Q.
\]
Poisson moments are bounded uniformly in $j$ and the series converges
for $Q>2k$. The same argument after differentiation of a slightly
broader Gaussian proves the gradient bound. The Fourier bound follows
from the heat factor and the characteristic function of a probability
measure.

In the coordinates moving with $2b$, write
$\partial_{t,b}=\partial_t+2b\cdot\nabla_x$, and also convolve in time
by a Gaussian of width
\[
 T=T(K)=KH/P_K,
 \qquad P_K=\min(H,K/H)^\theta,\qquad0<\theta<1/4,
\]
for $H\le K\le h$, using any increasing continuation below $H$.
Write $\rho_{K,b},F_{K,b}$ for these convolutions. The spatial
generator in logarithmic scale is, up to fixed heat conventions,
$L_K=K^2\Delta_x+j_K*-I$.

\smallskip\noindent\textbf{Generator bounds.}
We claim that for some $c'>0$
\begin{equation}\label{f:rho-generator}
 H^{-k}\sum_b\|L_K\rho_{K,b}\|_2^2
 \lesssim S^C\|\rho\|_1\min(H,K/H)^{-c'}.
\end{equation}
Indeed \eqref{f:ballcount} and the Gaussian identity for Fourier ball
mass give
\[
 \int_{|\zeta|\lesssim1/r}|\widehat\rho(\zeta)|^2\dd\zeta
 \lesssim S^C\|\rho\|_1(1+H^2/r),\qquad r\ge H.
\]
To see the normalization, the local mass is at most
$S^Cr^k(H^2+r)$ and the smoothing volume is $r^{k+1}$.
For $r<H$ the same bound loses only a fixed polynomial in $H/r$.
On $|\xi|\asymp1/r$ the squared spatial multiplier is bounded by
$C_N(K/r)^2(1+K/r)^{-N}$. For $r\le T$, the time Gaussian
restricts $|\tau+2b\cdot\xi|\lesssim T^{-1}$, with rapidly decreasing
dyadic tails. Of the $O(H^k)$ cap centers, a fraction at most
$C(r/T+H^{-1})$ obeys this restriction; this is the elementary count
of an $H^{-1}$ grid in a slab. The corresponding spacetime Fourier
ball has radius $O(r^{-1})$. For $r\ge T$ use radius $O(T^{-1})$ and
the trivial cap fraction one. At $r=K$ the resulting factor is
\[
 (P_K/H)(1+H^2/K)
 \le 2\min(H,K/H)^{\theta-1}.
\]
The dyadic sums for $r<K$, $K<r<T$, and $r\ge T$ converge, after
choosing $N$ large; enlargements in the Gaussian tails cost only fixed
powers. This proves \eqref{f:rho-generator}, also with the heat and
jump multipliers treated separately.

The exact equation supplies a second estimate:
\begin{equation}\label{f:current}
 |\partial_{t,b}F_{K,b}|\lesssim(HK)^{-1}F_{K,b}.
\end{equation}
After carrier removal $U_b$ has frequency support $O(H^{-1})$, and
its exact equation gives
\[
 \partial_{t,b}|U_b|^2
 =-2\nabla_x\!\cdot\operatorname{Im}
       \bigl(\overline{U_b}(\nabla_x-ib)U_b\bigr).
\]
Move this divergence to $J_K$. The required weighted Bernstein bound is
\[
 \int j_K(x-y)|(\nabla_y-ib)U_b(y,t)|^2\,\dd y
 \lesssim H^{-2}\int j_K(x-y)|U_b(y,t)|^2\,\dd y.
\]
To prove it, remove the carrier and reproduce with a smooth projection
kernel $R_H=H^{-k}R(\cdot/H)$. Polynomial comparison of translates of
$j_K$ and weighted Young's inequality give operator norm at most
\[
 \int|\nabla R_H(u)|\braket{u/K}^{Q/2}\,\dd u\lesssim H^{-1}
 \qquad(K\ge H).
\]
Together with \eqref{f:kernel} and Cauchy--Schwarz this bounds the
current by $C(HK)^{-1}(J_K*_xF_b)$. The sole time-weight term satisfies
$|w_h'|\le C_Jh^{-2}w_h$, and $h^{-2}\le(HK)^{-1}$.
Positive comoving time convolution preserves the inequality. The
argument requires only compact-frequency $L^2$ data for the summed
$U_b$; smooth approximation in that fixed cap justifies the current
identity and the limiting convolutions.

\smallskip\noindent\textbf{Propagation of the clipped bound.}
Put $\phi(u)=M_c(1-e^{-u/M_c})$ and
$I(K)=\sum_b\int\rho_{K,b}\phi(F_{K,b})$. For any symmetric
diffusion or jump generator $L$, concavity gives
$\phi'(F)LF\ge L\phi(F)$. Thus differentiating $I$ gives a lower
bound consisting of twice the generator pairing with $\rho$.
For the spatial part Cauchy--Schwarz, \eqref{f:rho-generator}, and
$\phi(F)^2\le M_cF$ yield
\[
 |\text{spatial error}|
 \le S^C\|\rho\|_1\min(H,K/H)^{-c'/2},
\]
where we used $D\le S^C\|\rho\|_1$ from \eqref{f:Sbalance}.
Notice the cancellation of the factors $H^{k/2}$ and $M_c^{1/2}$.
For the time part integrate by parts once. The Gaussian gives
$\|T\partial_{t,b}\rho_{K,b}\|_1\lesssim\|\rho\|_1$, while
\eqref{f:current} and $u\phi'(u)\le CM_c$ give cost
$CM_cT/(HK)$ per cap. There are $O(H^k)$ caps, so the time error is
at most $S^C\|\rho\|_1/P_K$. Both errors integrate from $K_0$ to $h$
to at most
\[
 S^C\|\rho\|_1\{S^{-cL}+H^{-c}\log h\}=o(\|\rho\|_1)
\]
when $L$ is fixed sufficiently large. We used $H\ge h^{2c/k}$ and
$S=h^{o(1)}$. Smooth regularization of $\rho$ and piecewise smooth
approximation at $K=H^2$ justify the differentiation; the positive
convolution formulas allow passage to the limit.

At $K=h$, drop clipping. Convolving twice averages $w_h|U_b|^2$ on
spatial width $h$ and time width at most $Hh$ along $2b$. The original
transition masks remain outside this estimate. For a witness
$z=(x_z,t_z)$, the Gram bound for the uncut original fine packets is
\begin{equation}\label{f:original-witness-Gram}
 \sum_b\int j_h(x-x_z-2b\tau)|U_b(x,t_z+\tau)|^2\,\dd x
 \lesssim\braket{\tau/(Hh)}^{M_1}\sum_Ts_T^2w_T(z).
\end{equation}
The relative displacement in width-$h$ units is $O(|\tau|/(Hh))$,
and elapsed native time is at most
$C+|\tau|/h^2\le C+|\tau|/(Hh)$. The prescribed finite localization
seminorms of the individual original packets grow polynomially in
this variable. Spatial almost orthogonality and summable Gram rows
therefore prove \eqref{f:original-witness-Gram}, retaining the original
exponent on $w_T$ by choosing the auxiliary kernel order $Q$ larger.
Since $w_h\le1$ and $T(h)\le Hh$, the Gaussian moments give
\[
 I(h)\lesssim\int\rho(z)\sum_Ts_T^2w_T(z)\lesssim\|\rho\|_1.
\]
At $K_0$, kernel comparability in space and \eqref{f:current} in
comoving time show that $F_{K_0,b}$ varies by bounded factors in
its central smoothing box. Therefore the extra convolution against
$\rho$ dominates a constant multiple of
$\sum_{z\in Y}H^{k+2}\sum_b\min(F_{K_0,b}(z),M_c)$.
We have proved that this last clipping sum has bounded mean on $Y$.

\smallskip\noindent\textbf{Recovery of packet envelopes.}
We now bound the actual square envelopes by the smoothed auxiliary
density on the test interval. This also fixes the required finite
orders. Partition the exact data of $U_b$
into smooth caps of width $K_0^{-1}$. Form the actual packets by
localizing $\chi_h P_vU_b$ in their scale-$K_0$ moving coordinates,
and retain the original $H$ masks externally. Since $K_0\le h$, fixed
native derivatives of $\chi_h$ are bounded and its temporal Fourier
thickness is at most $O(K_0^{-2})$. Use weighted suprema as packet
sizes. Local constancy, $|\chi_h|^2\le C_Jw_h$, and bounded overlap
of the spatial projections give
\begin{equation}\label{f:actual-packet-energy}
 K_0^{k+2}\sum_{v,\lambda}s_{b,v,\lambda}^2
 \lesssim\sum_v\int|\chi_hP_vU_b|^2
 \lesssim\int w_h|U_b|^2.
\end{equation}
Thus \eqref{f:auxiliary-energy} supplies the total integrated budget
$CD$ for this representation of the same actual cutoff field.

Here is an explicit fixed-order comparison. Work in native scale
$K_0=1$ for a cap, write its uncut carrier-removed solution as $q_v(u,s)$,
and take
\[
 N=Q+k+2,\qquad P_{\rm work}=3N+3.
\]
A smooth larger frequency projection has propagation kernel
$|L_\sigma(r)|\le C_N\braket{\sigma}^{N}\braket r^{-N}$.
Consequently, for every fixed $s_*$,
\[
 |q_v(u,s)|^2\le C\braket{s-s_*}^{2N}
        \int\braket{u-y}^{-N}|q_v(y,s_*)|^2\dd y.
\]
The bounded multiplier $\chi_h$ and its fixed native derivatives only
change the constant when estimating the actual packet sizes from this
uncut field. Choose the size-supremum exponent at least
$\max(P'_{\rm work},2N)$ for some fixed
$P'_{\rm work}>P_{\rm work}$. For the packet centered at
$(a,m)\in\Z^{k+1}$ its size then obeys
\[
 s_{v,a,m}^2\le C\braket{m-s_*}^{2N}
              \int\braket{a-y}^{-N}|q_v(y,s_*)|^2\dd y.
\]
Summing with the work packet weight of order $P_{\rm work}$ first in $m$
leaves time decay $P_{\rm work}-2N-1=N+2$, and then summing in $a$
retains spatial decay at least $N$. Thus the packet envelope at $z$
is bounded by a spatial $j_{K_0}$-weighted energy at any time
$|t_*-z_t|\le T(K_0)$. Here $T(K_0)\le K_0^2$ and the relative
velocity displacement is $O(T/H)\le CK_0$.

For these cap energies use weighted Bessel:
\[
 \sum_v\int j_{K_0}(x-x_0)|P_vf(x)|^2\dd x
       \lesssim\int j_{K_0}(x-x_0)|f(x)|^2\dd x.
\]
One proof avoids a nonsummable bound for each projection. With an
even polynomial weight, differentiate the cap symbols and use their
bounded overlap and Plancherel to bound synthesis
$(g_v)\mapsto\sum_vP_v^*g_v$ in positive polynomially weighted
$L^2(\ell^2)\to L^2$. Duality gives the displayed negative-weight
analysis estimate. A partition into boxes and summable convolution
extends it to translates of $j_{K_0}$. Apply it to $U_b$ at the fixed
time $t_*$. Average the preceding packet comparison over the central
part of the time Gaussian. For $z$ on the test slab and
$|t_*-z_t|\le T(K_0)\le h^2$, the polynomial weight $w_h(t_*)$ is
bounded below by a fixed constant. No lower bound for $\chi_h$ outside
the test slab is needed. Hence
$e_b^2(z)\lesssim F_{K_0,b}(z)$ on the test interval. Markov's
inequality, with any slowly diverging subpower threshold, now proves
\eqref{f:clipped}.

When $K_0\le h$, these reconstructed packets are the leaves of the
subpower tree: no original scale-$h$ fine squares are inserted below
them. Their rapid output decay can be fixed above $P'_{\rm work}$,
and the same work exponent is used in each later application of
\Cref{p:weighted,p:hierarchy,p:near-cauchy}. The original order $P$
enters only the input square hypothesis and the original-packet Gram
bound used at $K=h$; no comparison of a reconstructed fine array
with the original square envelope is asserted. If $K_0>h$, use the
unchanged original scale-$h$ representation as the leaf, with order
$P$. Its square envelope is bounded at the witnesses, so clipping is
automatic. All orders are chosen before the finite scale grids and
the finite original derivative order in \Cref{f:exact}; they depend
on $k,c$ and the original envelope order, not on hierarchy depth.
\end{proof}

Retain only original $H$ masks whose centers obey $e_b^2\le M_c$.
The loss is negligible. Indeed, within a sufficiently small fixed
power enlargement, a removed bin has $e_b^2\gtrsim M_c/S$ at the
witness, by polynomial weight comparison. Formula~\eqref{f:clipped}
therefore bounds the number of such bins by
$H^kS^{1-C_3+o(1)}$. Multiply this by the original amplitude
$SH^{-k/2}$ and choose $C_3$ large. Outside the enlargement, use the
rapid cutoff tails and the same absolute amplitude bound.

At surviving transition masks use weights $b_T=e_b(z_T)$. For each
bin $b$, sampling only the retained mask centers gives
\[
 \sum_{T\text{ retained in }b}e_b^2(z_T)w_T(z)
 \lesssim \min\{e_b^2(z),M_c\}.
\]
Convolution gives the bound by $Ce_b^2(z)$; the center restriction
$e_b^2(z_T)\le M_c$ and the bounded sum of the mask weights give the
bound by $CM_c$. At strict ancestors we may therefore use the square
sum of $\min(e_b^2,M_c)$ over descendant bins, with proportional
clipping inside each bin. At finer scales, up to the chosen leaf
$K_0$ or $h$, use its full descendant squares with the same envelope
order. The displayed comparison verifies \Cref{p:hierarchy} at the
crossing of the clipping threshold; ordinary convolution verifies the
other transitions. Root and transition square sums are $S^{o(1)}$.
Sorting actual transition amplitudes by their ratios $r$ to $b_T$,
Proposition~\ref{p:weighted} from the endpoint gives energy
$S^{o(1)}r^{2-q}D$; the square bound of this class is $S^{o(1)}r$.
The ratios are polynomially bounded in $S$ by Cauchy, and negligible
negative powers make no root contribution. \Cref{p:near-cauchy}, with its exponent slack sent to zero, therefore
gives
\begin{equation}\label{f:subbalance}
 |Y|\le S^{o(1)}D_iH_i^{-k-2},\qquad
 H_i/H_g=S^{o(1)},\qquad |Y|H_i^{k+2}/D_i=S^{o(1)}.
\end{equation}
The reverse comparisons follow from \eqref{f:selected-counter}.
Rounding $H_i$ or the fine endpoint by fixed factors has no effect;
if $H_i\gtrsim h$, use the original $h$ label grid.

\subsection{Finite scale trees and the limiting profiles}

The logarithmic argument will always be realized on a fixed finite
grid before limits are taken. In the power case use
\[
 \kappa_i(s)=h^s,\quad0<s<e_i=1,\quad\Lambda=h.
\]
In the subpower case use
\[
 \kappa_i(s)=H_iS^s,\quad-\infty<s<e_i,
 \quad e_i=\min\{L,\lim\log_S(h/H_i)\},\quad\Lambda=S.
\]
The actual fine endpoint and the original transition masks are
included in every tree. A transition mask is a distinct linear layer
even if its exponent coincides with the endpoint exponent. Physical
lattices need not be nested. Frequency labels retain their assigned
ancestors despite bounded overlaps of the enlarged supports.

We spell out the selection order because subsequent directional
deletions change intermediate amplitudes. No assertion below assumes
that a selected amplitude ratio survives such a deletion.

\begin{lemma}[Finite-grid selection and discard]\label{f:selection}
On a fixed finite grid one can select ratios $r_{i,s}$, retain a
subpower fraction of the high cubes, and retain packet masks with the
following properties. Put
\begin{equation}\label{f:occupancy-size}
 M_{i,s}=\kappa_i(s)^{k+2}(r_{i,s}/A_i)^q.
\end{equation}
At every selected scale the original selected amplitudes obey
$a_T\asymp r_{i,s}b_T$, and
\begin{align}
 Q_{i,s}&:=\sum_T|T|a_T^qb_T^{2-q}=\Lambda^{o(1)}D_i,
                                                        \label{f:Qsat}\\
 E_{i,s}&:=\sum_T|T|a_T^2=\Lambda^{o(1)}D_ir_{i,s}^{2-q},
                                                        \label{f:Esat}\\
 \sum_{T\text{ active}}b_T^2&\le
             \Lambda^{o(1)}|Y|/M_{i,s}.                  \label{f:Bcharge}
\end{align}
An active cell has at most $\Lambda^{o(1)}M_{i,s}$ eligible witnesses
in its prescribed subpower enlargement. All witnesses are eligible in
the power case and at strict ancestors of transition bins. Within a
subpower transition bin $b$, eligibility means
$e_b^2(z)\le\Lambda^{o(1)}M_c$.

More generally, a collection of actual outputs in fixed masks whose
weighted $q$ budget is at most $\Lambda^{-\eps}D_i$, for fixed
$\eps>0$, has negligible root contribution. The same assertion holds
for error summands in any later restricted tree.
\end{lemma}

\begin{proof}
Fix a small number $\beta>0$ which bounds the exponents of all norm
losses, fixed-grid constants, logarithmic sums, and initial balance
errors. Work far enough along the sequence that these losses are at
most $\Lambda^{C\beta}$. Choose the distance cutoff exponent next,
the eligibility and enlargement exponents after that, and the
occupancy cutoff exponent last. Each is larger than the fixed
comparability multiples of its predecessors. They will all tend to
zero, in this order, along a slow diagonal.

First prove the discard statement, without assuming selection or
saturation. For an error $U$ with actual amplitudes and budget $Q_U$,
Propositions~\ref{p:weighted} and \ref{p:hierarchy} through its masks
give, in the power case,
\begin{equation}\label{f:discard}
 \#\{z\in Y:|U(z)|>tA_i\}
 \le\Lambda^{C\beta}t^{-q}A_i^{-q}Q_U.
\end{equation}
One can jump directly from a current layer to the root: no not-yet
assembled intermediate masks are needed, and any transition
restriction already fixed remains imposed. In the subpower case
propagate to the gap edge. Sort the \emph{actual} output amplitudes
there by their ratios to the propagated weights. Cauchy bounds the
nonnegligible ratios by fixed powers of $S$, since only finitely many
scales are involved and the clipped square is bounded at the
transition. \Cref{p:near-cauchy} across the gap gives
\eqref{f:discard} with a further $S^{C\eta}$, where $\eta>0$ is its
arbitrarily small exponent slack. Choose $\eta$ after the desired
discard saving. Very small ratios have negligible contribution by
absolute Cauchy, and larger ratios have summably small weighted
budget. Thus $Q_U\le\Lambda^{-\eps}D_i$ allows both a vanishing
threshold fraction and a power-small exceptional set, by choosing
$\beta,\eta$ and the exponent in $t$ sufficiently small relative to
$\eps$. This proof uses actual error amplitudes throughout.

Assemble from leaves toward the root, with the original transition
masks fixed. At each stage Cauchy gives a polynomial upper bound for
the amplitude-to-weight ratios. In the subpower case this includes
the quotient $\kappa_i(s)^{k/2}/A_i$, which is a fixed power of $S$ at
each fixed scale. Ratios below a sufficiently small inverse power
are harmless by Cauchy and the discard estimate. There remain
$O(\log\Lambda)$ dyadic classes. Select a class for the sum actually
reaching the root before packetizing a smaller inverse width. A finite
number of selections costs $\Lambda^{o(1)}$ in thresholds and high
set size.

Here is the occupancy deletion at this same stage. If $n_T$ counts
eligible witnesses in the prescribed enlarged cell, local weight
comparison and \eqref{f:clipped} give
\[
 \sum_Tb_T^2n_T\le\Lambda^{\beta'}|Y|,
\]
where $\beta'$ accounts for the already chosen enlargement and
eligibility. This also holds with full squares in the power case.
For a candidate ratio class $r$ delete cells with
$n_T>\Lambda^\gamma M_s$, choosing $\gamma>\beta'+C\beta$. Their
weighted charge is
\[
 Q_{\rm bad}\lesssim r^q\kappa_s^{k+2}
                  \sum_{\rm bad}b_T^2
 \le\Lambda^{\beta'-\gamma}|Y|A_i^q
 \le\Lambda^{\beta'-\gamma+C\beta}D_i.
\]
Use \eqref{f:hbalance} or \eqref{f:subbalance} in the last inequality.
The discard estimate therefore permits this pruning before the final
selection at that layer. Similarly, a size supremum realized only at
distance $\Lambda^\gamma$ from the packet center carries a negative
power budget if its local amplitude is power-smaller: retain an extra
fixed decay margin in the supremum and apply the same estimate.

The upper bound for $Q_{i,s}$ follows from the weighted norm estimate
through the already assembled masks. The reverse follows by applying
\eqref{f:discard} to the entire selected field, which still has its
root high set. Since $a_T\asymp r_{i,s}b_T$, the formulas for $E_{i,s}$
and $\sum b_T^2$ follow algebraically. Taking the diagonal of the
specified slacks proves the displayed subpower assertions.
\end{proof}

\begin{lemma}[Profiles]\label{f:profiles}
After taking subsequences and finer finite rational grids there are
nonnegative, nonincreasing, $k/2$-Lipschitz functions $f_i$ such that
\[
 r_{i,s}=\Lambda^{f_i(s)+o(1)},\qquad f_i(e_i)=0.
\]
In the power case $f_i(0)=\lim\log_h A_i$. In the subpower case
\begin{equation}\label{f:d-bounded}
 0\le d_i(s):=f_i(s)+ks/2\le ke_i/2\qquad(s\le0).
\end{equation}
\end{lemma}

\begin{proof}
For $s<s'$, weighted Cauchy between the two layers gives
\[
 r_{i,s}\le\Lambda^{o(1)}
             (\kappa_i(s')/\kappa_i(s))^{k/2}r_{i,s'}.
\]
Conversely, integrated spatial $L^2$ orthogonality gives
$E_{i,s}\le C_{\rm grid}E_{i,s'}$. To justify it with masks, first
remove the outermost mask by its bounded sup norm; then separate its
immediate children by their bounded Fourier overlap. Repeat in this
order. One never commutes all outer masks past a final fine frequency
decomposition. Polynomial spatial weights can be retained using
translated band-limited cutoffs. Formula~\eqref{f:Esat}, and $q>2$,
now give $r_{i,s}\ge\Lambda^{-o(1)}r_{i,s'}$. At the fine endpoint
the ratio is one by the choice of the input sizes.

In the power case root witness sampling and local constancy give
$|Y|A_i^2\lesssim\Lambda^{o(1)}E_{i,s}$. Together with saturation
and $|Y|=\Lambda^{o(1)}D_iA_i^{-q}$ this gives
$r_{i,s}\le\Lambda^{o(1)}A_i$. The reverse endpoint comparison
comes from root Cauchy,
$A_i\le\Lambda^{o(1)}\kappa_i(s)^{k/2}r_{i,s}$, as $s\downarrow0$.
The same root Cauchy in the subpower case gives
$f_i(s)+ks/2\ge0$. The fine endpoint and the Lipschitz bound give
the other inequality in \eqref{f:d-bounded}.

On each finite grid the logarithms are uniformly bounded and have
these two one-sided Lipschitz bounds up to vanishing errors. A
diagonal subsequence on rational parameters gives the asserted
continuous profiles. Every later test involves only a fixed finite
grid: include its parameters first and take the subsequence far
enough that its logarithms approximate the limiting values to the
desired accuracy. Grid complexity can equivalently increase
sufficiently slowly. No rate of convergence to a particular slope is
assumed.
\end{proof}

Freeze the selected masks together with the ratios and budgets recorded
in Lemma~\ref{f:selection}. These records will charge later deletions
and supply the broad witnesses. The original transition amplitudes in
\eqref{f:original-amplitudes} remain separate records for the final
incidence count. After a deletion the fields receive fresh actual
amplitudes and fresh propagated positive weights; neither set of
recorded amplitudes is asserted to survive that deletion.

\subsection{Positive paths and the occupancy identities}

We next turn the selected packet hierarchy into a law on directions
through a high point. Its purpose is to compare two quantities: the
square mass available in a frequency cap, and the number of high points
in a physical cell directed by a sampled leaf. The normalization of the
law must be recomputed when directions are removed, although its
occupancy scale $M_{i,s}$ and its charging bounds come from the frozen
selection.

Fix one field and one of its finite trees. Initially use the selected
tree itself. Later we use the same construction after the particular
angular omissions of Lemma~\ref{f:angular}. Those omissions keep the
frozen masks and retain the root high threshold. Their fresh weighted
norm bound is proved in that lemma using the packet estimates and the
discard statement; the construction below uses that bound, not any
assertion that the old amplitudes survive the omissions.

Expand the chosen positive square propagation through the retained
cells, keeping their assigned cap ancestry. A surviving path
$\pi=(T_0,\ldots,T_\ell)$ ends at a terminal physical input
$\gamma=T_\ell$. Denote its successive nonnegative propagation factors
by $k_0(z;\pi),\ldots,k_\ell(z;\pi)$. The first evaluates at the root
point $z$; an ordinary parent-to-child factor is the native sampling
weight $w_{T'}(z_T)$ from \eqref{p:parent-weights}, with its chosen
fixed normalization and edge-decay reserve. Retain each factor's
dependence on the assigned cap ancestry and its evaluation point.
With the input square $s_\gamma^2$, the weight of the path is
\[
 s_\gamma^2\prod_{j=0}^{\ell}k_j(z;\pi).
\]
Let $\omega(\gamma)$ denote its terminal direction label. Sum these
weights over the surviving paths with $\omega(\gamma)=\omega$ to define
$J_z(\omega)$, and put
\[
 W(z)=\sum_\omega J_z(\omega).
\]
In particular, terminal physical inputs with the same direction label
are summed. Wherever $W(z)>0$, the provisional direction law is
$p_z(\omega)=J_z(\omega)/W(z)$. The proof removes distant paths and
normalizes once more; all conclusions below concern that resulting
near-path law. Later angular omissions likewise require fresh path
sums and a fresh normalization.

The reference cap weights $b_\tau^2(z)$ retain the available descendant
squares of the selected tree. In the power case these are full squares.
In the subpower case strict transition ancestors use proportionally
clipped bins, and finer levels use full descendant squares only in
eligible transition bins. This retains the transition-center restriction
and the eligibility condition in Lemma~\ref{f:selection}. Indeed,
in the subpower case,
\[
 \sum_{b\text{ eligible}}e_b^2(z)
 \leq\Lambda^{o(1)}\sum_b\min(e_b^2(z),M_c)
 \leq\Lambda^{o(1)},
\]
because eligibility gives $e_b^2(z)\leq\Lambda^{o(1)}M_c$.
Thus these level-specific reference squares remain controlled even
where full squares are used. They are distinct from the surviving
path mass $J_z$. Their comparison with the law will hold at sampled
labels, not uniformly at every label after omissions.

For finite-valued labels $A$ and $B$, $\Pr(A\mid B)$ denotes
$p_{A\mid B}(a\mid b)$ evaluated at the sampled values $a=A$, $b=B$.
Its logarithmic information in base $\Lambda>1$ is
$-\log_\Lambda\Pr(A\mid B)$. In
$\Pr(T_s(Z,\Omega_i)\mid\Omega_i)$, conditioning on $\Omega_i$ fixes
the tilted lattice, and $T_s(Z,\Omega_i)$ is its sampled cell label.

\begin{lemma}[Restricted path law]\label{f:paths}
For the selected tree, and for each of the directional restrictions
constructed below, one can draw a uniform root point $Z$ from a
subpower fraction of $Y$ and then a path to a fine direction $\Omega_i$
in each field. The intermediate active cells lie within subpower
normalized distance of $Z$. The fields' paths may be drawn
independently conditional on $Z$.

For the reference cap square weights just defined, at a typical
sampled label
\begin{equation}\label{f:cap-law}
 \Pr(\tau\mid Z)=\Lambda^{o(1)}b_\tau^2(Z).
\end{equation}
Let $T_s(Z,\Omega_i)$ be the cell containing $Z$ in the disjoint
scale-$\kappa_i(s)$ lattice tilted by $\Omega_i$. At a typical draw,
\begin{equation}\label{f:occupancy-law}
 \frac{\Pr(T_s(Z,\Omega_i)\mid\Omega_i)}
      {\Pr(Z\mid\Omega_i)}=\Lambda^{o(1)}M_{i,s}.
\end{equation}
Typical means outside probability tending to zero on each fixed
finite test, with any prescribed positive exponent tolerance.
\end{lemma}

\begin{proof}
We first show that the available high points have enough surviving
path mass to normalize the law.
The weighted norm estimate applied afresh to this restricted tree
gives a budget at most $\Lambda^{o(1)}D_i$, with its actual modified
amplitudes. At the root its integrand is
\[
 |V(z)|^q W(z)^{1-q/2}.
\]
The number of high points where $W\le\Lambda^{-\alpha}$ is at most
$\Lambda^{o(1)-\alpha(q-2)/2}D_iA_i^{-q}$. Across the subpower gap
the ratio-class argument in the proof of Lemma~\ref{f:selection}
adds only $S^{C\eta}$; choose $\eta$ smaller than
$\alpha(q-2)/(4C)$. Thus low path mass is exceptional. Full-background
convolution and the transition eligibility test give
$W\le\Lambda^{o(1)}$.

Use a fixed extra polynomial decay power on each path edge. A path
with an intermediate center farther than $\Lambda^\alpha$ from $z$
has a comparably large normalized edge displacement, since the
inverse widths increase toward the leaves and the nested-cap shears
are bounded in the appropriate coordinates. Dropping the extra
power on this edge still leaves a convolution bounded by the full
background; hence such paths have mass $o(W(z))$ after the fixed
extra margin is chosen larger than the path-mass lower-bound
exponent requires. This uses the same fixed margin on each edge,
not an increasing endpoint weight exponent. Redefine $J_z$ and $W$
to count only the retained near paths, and normalize by this new
$W$. Relabel the final retained root set as $Y$; its cardinality
differs by only a subpower factor from the preceding set. Every
retained path is eligible at its transition bin: its surviving
transition-mask center has $e_b^2\le M_c$, and proximity to that
center and polynomial weight comparison give
$e_b^2(z)\le\Lambda^{o(1)}M_c$. On the near paths, telescoping local
weight comparisons and the subpower neighboring-cell multiplicities
give
\[
 p_z(\omega):=J_z(\omega)/W(z)
       \le\Lambda^{o(1)}a_\omega(z),
\]
where $a_\omega$ is the corresponding available leaf square weight.
If a final direction label sums many final input cells, only its
last localization layer and intermediate active cells need be near
$z$; the input coefficients are summed in $a_\omega$.

Summing this upper bound over descendants proves the upper half of
\eqref{f:cap-law}. For its lower half, at any fixed root point the
total probability of sampled labels satisfying
$p_z(\tau)<\Lambda^{-\eps}b_\tau^2(z)$ is at most
$\Lambda^{-\eps}\sum_\tau b_\tau^2(z)=\Lambda^{-\eps+o(1)}$.
This is a statement at sampled labels, not a uniform lower bound
for all labels.

Here are both sides of \eqref{f:occupancy-law}. For a fixed leaf
direction $\omega$ and a tilted $s$ cell $C$, put
\[
 \mu(\omega,C)=|Y|^{-1}\sum_{z\in Y\cap C}p_z(\omega).
\]
At the realized pair $(z,\omega)$ the ratio in question is exactly
\begin{equation}\label{f:fiberidentity}
 |Y|\mu(\omega,C)/p_z(\omega).
\end{equation}
For the upper bound, all contributing active $s$ cells lie in a
subpower enlargement of $C$. The occupancy deletion counts at most
$\Lambda^{o(1)}M_{i,s}$ eligible points there. The leaf background
varies by a subpower factor within the tilted cell, and path
domination and the typical leaf lower bound compare each summand
to $p_z(\omega)$. This proves the upper bound.

For the lower bound, let $J_T^{\rm dom}(\omega)$ be the
unnormalized sum of positive descendant-path weights from an active
scale-$s$ mask $T$, evaluated at its center $z_T$ on the selected tree
before angular deletions.
Retain the transition-center restriction and proportional clipping,
and sum all terminal physical cells with direction label $\omega$.
Later macro-dependent omissions give subcollections of this one
majorant, without creating extra copies. Positive convolution gives
$\sum_\omega J_T^{\rm dom}(\omega)\lesssim C_{\rm grid}b_T^2$.
At the transition this uses both the full-square convolution bound
and the bound by $C_{\rm grid}M_c$ from retained centers; take their
minimum. At finer cells use the full descendant weights.

Let $\mathcal N_s(\omega,C)$ consist of active masks in the assigned
$s$-cap whose $\Lambda^\eta$ native enlargements meet the tilted
$s$-cell $C$, and put
\[
 a_{\omega,C}=\sum_{T\in\mathcal N_s(\omega,C)}J_T^{\rm dom}(\omega).
\]
This reference charge is used only for domination; the actual law
$p_z(\omega)=J_z(\omega)/W(z)$ is normalized using the fresh surviving
mass of the restricted tree.
The near-ancestor kernel sum and $W(z)\ge\Lambda^{-o(1)}$ give the
following pointwise bound. For fixed $(\omega,T)$, only
$\Lambda^{o(1)}$ cells $C$ occur, because the two scale-$s$ frames
differ by a bounded normalized shear. Consequently
\[
 p_z(\omega)\le\Lambda^{o(1)}a_{\omega,C},\qquad
 \sum_{\omega,C}a_{\omega,C}
 \le\Lambda^{o(1)}\sum_{T\text{ active}}b_T^2
 \le\Lambda^{o(1)}|Y|/M_{i,s}.
\]
At a strict transition ancestor, eligible contributions are comparable
at the charging center up to the prescribed subpower factor; unique
label ancestry charges each clipped bin only once.
The total probability of fibers satisfying
\[
 \mu(\omega,C)<|Y|^{-1}\Lambda^{-\eps}M_{i,s}a_{\omega,C}
\]
is at most $\Lambda^{-\eps+o(1)}$. For every other sampled fiber,
\eqref{f:fiberidentity} proves the lower bound. Replacing boxes by
balls or bounded-overlap coverings changes these sampled logarithms
by $o(1)$: the total probability of a cell power-smaller than one
of its boundedly many neighbors is power-small. The same argument
handles subpower distortions.
\end{proof}

\begin{lemma}[Conditioning directions]\label{f:conditioning}
Inside a fixed cap $\tau_l$, let $C$ be a physical cell over which
every descendant velocity has drift at most the endpoint width, and
whose time range is at most the endpoint duration. Adding an
independently redrawn descendant direction as conditioning changes
the logarithmic information of position inside $C$ by $o(1)$ at a
typical draw. The cell may depend on a different conditionally
independent drawn direction. The conclusion also holds for
subpartitions and subpower distortions. In the power case it holds
for a direction from another field whenever its full weights have
the same bounded drift on $C$.
\end{lemma}

\begin{proof}
Use as reference direction law the endpoint square weights at the
center of $C$, normalized within $\tau_l$. Clip whole transition
bins proportionally if $\tau_l$ is a strict ancestor. The drift
assumption makes these reference weights comparable throughout $C$.
By \eqref{f:cap-law} and path domination, the conditional density of
the direction relative to this law is at most $\Lambda^{o(1)}$ at
typical data. Its average conditional density is at least
$\Lambda^{-o(1)}$ at a typical sampled direction: the exceptional
mass is bounded by summing reference weights exactly as in the
lower half of \eqref{f:cap-law}. Bayes' identity now bounds the
likelihood ratio for position by $\Lambda^{o(1)}$.

For clarity, the remaining one-sided information comparison requires
no independence assumption. If $p$ and $p'$ are the probabilities
of a realized finite label before and after additional conditioning,
then $\mathbb E[p/p']\le1$. Markov's inequality implies
$p/p'\le\Lambda^\eps$ outside probability $\Lambda^{-\eps}$.
Apply this to each member of a finite chain of partitions. Combining
it with the reference-law comparison proves equality of the relevant
information increments to arbitrary logarithmic tolerance.
\end{proof}

\subsection{Angular pruning inside one profile block}

\begin{lemma}[Angular test and original broad witnesses]\label{f:angular}
Fix one field, a cap parameter $l$ away from the endpoints and the
transition seam, and a finite mesh of $0<v\le u$ with $l+Cu<e_i$.
Suppose that
\[
 f_i(l)-f_i(l+v)\ge xv
\]
on this mesh, for some $x>0$. For a number $\eta>0$ depending only on
$x,k$, one may restrict the frozen tree so that, at typical $(Z,\tau_l)$,
its new conditional direction law assigns mass at most
$\Lambda^{-\eta v+o(1)}$ to every affine hyperplane strip of width
$\Lambda^{-v}$ in the native frequency coordinates of $\tau_l$.

At one arbitrarily small mesh precision $v_0>0$, the sampled active
$l$ cells may also be required to have broad witnesses for their
\emph{original pre-modification} fields. Namely, there are $k+1$
$(l+v_0)$ cap groups, transverse throughout their supports with
determinant at least $\Lambda^{-C_kv_0}$, each having magnitude at least
$\Lambda^{-C_kv_0}r_{i,l}b_{\tau_l}(Z)$ at a point within native
distance $\Lambda^{C_kv_0}$ of the cell.

The high root set loses only a vanishing fraction, and Lemmas
\ref{f:paths} and \ref{f:conditioning} apply to the modified tree with
the same recorded occupancy sizes. This construction is made for one
fixed block at a time; no simultaneous assertion over all blocks is
needed.
\end{lemma}

\begin{proof}
Suppress the field index. In the native coordinates of $\tau_l$, use
macro boxes of side $\Lambda^{3u}$. The endpoint descendant weights,
divided by the weight of $\tau_l$, vary within bounded factors on
these boxes: all affected cell widths and durations are at most
$\Lambda^{2u+o(1)}$, and the endpoint is separated by $Cu$. If a
transition ancestor occurs, clip each entire bin proportionally.

For each mesh value $v$, greedily remove all $(l+v)$ caps intersecting
an enlarged hyperplane neighborhood whose reference measure is at
least $\Lambda^{-2\eta v}$. There are at most
$O(\Lambda^{2\eta v})$ choices, and the residue has the desired upper
bound before normalization. Carry out the omissions from coarser
to finer mesh values, so earlier omissions only restrict subsequent
masks.

We verify that their contribution at level $l$ is cheap. Let
$q'=2(k+1)/(k-1)>q$, the lower-dimensional exponent. The original
selected budget at $l+v$ is bounded by
\[
 Q_{q',l+v}\le\Lambda^{o(1)}r_{l+v}^{q'-q}D.
\]
This follows from the recorded ratio class, before any modification.
The plate-supported version of Proposition~\ref{p:weighted} at
exponent $q'$ propagates each omitted hyperplane contribution through
\emph{all} intervening masks. Plate support is inherited by each
descendant in its own normalized coordinates. Summing the omitted
planes costs at most $\Lambda^{2\eta q'v}$. On the active $l$ masks,
H\"older with measure $b_T^2|T|$ and
$\sum b_T^2|T|\le\Lambda^{o(1)}Dr_l^{-q}$ gives
\begin{align*}
 Q_{q,l}(\text{error})
 &\le \big(\Lambda^{2\eta q'v+o(1)}
                 r_{l+v}^{q'-q}D\big)^{q/q'}
          \big(\Lambda^{o(1)}Dr_l^{-q}\big)^{1-q/q'}\\
 &\le D\Lambda^{2\eta qv+o(1)}
       (r_{l+v}/r_l)^{q(q'-q)/q'}.
\end{align*}
Choose $2\eta<x(q'-q)/(2q')$. The profile assumption then makes this
a negative fixed power budget, and Lemma~\ref{f:selection} discards
it at the root. Earlier restrictions are harmless: use actual
modified amplitudes when propagating, and bill the initial
restricted subfields to their original fine budgets.

We include the localization justification for the fact that the
direction subsets can vary by macro box. Assign each $l$ packet by
its center to such a box and alter the field \emph{before} its
Schwartz packet cutoff. No rough physical cutoff is inserted.
At a single hierarchy step, packets centered away from a target
macro box have an arbitrarily large fixed inverse polynomial
box-distance gain, by their cutoff and supremum tails. The
denominator weights may be left in the child coordinates, using
\Cref{p:hierarchy}; one need not compare them across an entire
macro box. These gains form a summable convolution in macro offsets.
An equivalent algebraic formulation is useful. For a summable
polynomial macro weight $\Psi$ and exponent $p>2$, replace
\[
 b_T\quad\hbox{by}\quad b_T\Psi(z_T)^{-1/(p-2)}.
\]
Then $a_T^pb_T^{2-p}$ acquires exactly the factor $\Psi(z_T)$.
Polynomial comparability of $\Psi$ requires only a fixed additional
edge-decay margin, since all affected widths are below the macro
scale. Use this for $p=q,q'$ at each step; the same margin suffices
at every depth. Summing the target-box estimates bills each initial
packet with bounded total weight. This proves the stated norm
charge for varying direction restrictions, and also proves the
fresh restricted-weight norm bound used in Lemma~\ref{f:paths}.

For broad witnesses, use separately macro boxes of side
$\Lambda^{3v_0}$ \emph{before} modifying directions. In each box
mark a group at level $l+v_0$ if its original contribution through
the frozen intervening masks has supremum at least
$\Lambda^{-C_kv_0}r_lb_{\tau_l}$ in a fixed enlargement. There are
at most $\Lambda^{kv_0+o(1)}$ groups, so for $C_k$ sufficiently large
the unmarked sum cannot supply a high witness defining an active
$l$ packet. The far-supremum deletion in Lemma~\ref{f:selection}
ensures such a witness occurs within the stated enlargement.

The elementary maximal-simplex test gives a dichotomy for the marked
centers: either $k+1$ of them have successive affine widths at least
a constant times $\Lambda^{-v_0}$, or all marked centers lie in a
constant enlargement of a hyperplane strip of that width. In the
first case the determinant is at least a fixed power of
$\Lambda^{-v_0}$; enlarge the separation constants so it holds
uniformly on the chosen supports. In the second case the same
$q'$ calculation shows that the marked contribution is cheap at
level $l$. Thus the nonbroad masks themselves have negligible
weighted $q$ budget and may be deleted first.

Finally reconstruct the positive path law for the modified fields,
as in Lemma~\ref{f:paths}. Low total path mass is exceptional by its
fresh norm bound. Its conditional law is dominated by the surviving
reference descendant weights divided by $b_{\tau_l}^2(Z)$; the
typical reverse cap comparison loses only a subpower. The heavy
strip cutoff $\Lambda^{-2\eta v}$ therefore implies the claimed
$\Lambda^{-\eta v+o(1)}$ angular bound. Occupancy upper bounds still
refer to the same frozen masks, and the charging proof of
\eqref{f:occupancy-law} is unchanged. The broad witnesses used later
are the recorded original subfields, whose selected amplitudes
remain available for norm and square estimates.
\end{proof}

\subsection{Ball inflation and the two allowed slopes}

\begin{lemma}[Ball inflation in a cap]\label{f:inflation}
Work inside a cap $\tau_l$ of the preceding lemma, and suppose
$f(l+v)-f(l)=-xv+O(\eps u)$ on the required finite mesh in $[0,2u]$.
Put $d_0=1-2x/k$. For the isotropic cell $Q_v$ of radius
$\Lambda^v$ in this cap's native coordinates, define the sampled
effective occupancy
\[
 N_v=\Pr(Q_v(Z)\mid\tau_l)/\Pr(Z\mid\tau_l).
\]
For $u\le v\le2u$ on that mesh,
\begin{equation}\label{f:inflation-occupancy}
 \log_\Lambda(N_v/M_l)\le(k+1)d_0v+O(\eps u).
\end{equation}
The error can be made arbitrarily small by taking the broad precision
$v_0/u$ small, refining the finite mesh, and then taking the sequence
limit. All spatial tests stay on one side of the transition seam.
\end{lemma}

\begin{proof}
Every point under the cap conditioning is near an active $l$ cell.
Such a cell has at most $\Lambda^{o(1)}M_l$ eligible points.
Path domination, the typical reverse cap comparison, and background
variation bound their relative conditional weights by subpowers.
It therefore suffices to count active $l$ cells in an isotropic
ball. They have the original broad witnesses from Lemma
\ref{f:angular}. Pigeonholing the separated tuple costs
$\Lambda^{O_k(v_0)}$, since there are only a fixed power of
$\Lambda^{v_0}$ possible labels.

Write $\kappa_s=\Lambda^s$ in the $l$ coordinates and, for the
$j$th broad group, let $\mathcal E_{j,s}$ be the local square sum of the original
selected packet amplitudes with their native weights at relative
scale $s$. The witness inequalities, including their allowed
offsets, imply at every counted unit cell
\[
 r_lb_{\tau_l}(Z)
 \lesssim\Lambda^{O_k(v_0)+o(1)}
             \prod_{j=1}^{k+1}\mathcal E_{j,v_0}^{1/(2(k+1))}.
\]
Masks between $l$ and $l+v_0$ have bounded absolute sums, and the
native weights absorb the witness offsets. The background
$b_{\tau_l}$ is comparable over the larger test ball.

For $s<t\le2s$, average the resulting product of squares with
exponents $1/k$ on a $\kappa_t$ ball. Proposition~\ref{p:kakeya}
gives
\begin{equation}\label{f:one-inflation}
 \frac1{|B_{\kappa_t}|}\int_{B_{\kappa_t}}
             \prod_{j=1}^{k+1}\mathcal E_{j,s}^{1/k}
 \lesssim\Lambda^{O_k(v_0)}
       \prod_{j=1}^{k+1}
          \left(\frac1{|CB_{\kappa_t}|}
                       \int_{CB_{\kappa_t}}\mathcal E_{j,s}\right)^{1/k}.
\end{equation}
Indeed these are nonnegative sums essentially constant on tubes of
width $\kappa_s$ and duration $\kappa_s^2\ge\kappa_t$. On the ball
replace them by columns crossing the ball, where each envelope is
comparable. Such a column costs a cross section of width
$\kappa_s$ times length at least $\kappa_t$ in the enlarged-ball
average. The quantitative determinant factor is
$\Lambda^{O_k(v_0)}$. Weighted offsets are summed with the fixed
integrable envelope powers.

We verify that each last average is bounded by the corresponding
$\mathcal E_{j,t}$, even when intermediate mask levels are skipped. This
verification prevents an aspect-ratio loss. Put $S_0=\kappa_s$ and
$T_0=\kappa_t$. Reproduction in the native $s$ coordinates bounds
the square of a weighted local supremum by a normalized weighted
$L^2$ integral of the pre-cutoff field. After summing $s$ cells,
the displacements from an averaged point have the form
\[
 (S_0U+2v_sS_0^2w,S_0^2w),
 \qquad\text{density}\ \le C(1+|U|+|w|)^{-P}.
\]
For each fixed displacement and time, integrate in space over width
$T_0$, replacing the ball by a smooth polynomial weight. Remove the
next mask by its bounded sup norm and separate its immediate child
caps by spatial almost orthogonality; localization at width $T_0$
preserves bounded Fourier overlap. Iterate in the actual order of
the masks. Translated band-limited cutoffs positive on boxes prove
the same assertion for the polynomial weight. At the final layer,
pointwise Cauchy bounds the fields by the $t$ packet envelopes.

In the frame of a final $t$ cap, the normalized shift just displayed
has sizes
\[
 (S_0/T_0)(U+O(1)w),\qquad(S_0/T_0)^2w,
\]
because its velocity differs from $v_s$ by $O(S_0^{-1})$.
Convolution with this contracted shift density preserves the order
$P$ of the final weight by the distance and volume comparison in
\Cref{p:hierarchy}. The auxiliary spatial ball weight has all
the finitely many moments needed for this comparison. Thus no
factor depending on $T_0/S_0$ is lost. One can use the same fixed
$P$ at each inflation step, with a fixed excess supremum margin.

At the terminal scale, the recorded ratio class and weight
convolution give
\[
 \mathcal E_{j,v}\lesssim\Lambda^{o(1)}r_{l+v}^2b_{\tau_l}^2(Z).
\]
Iterate \eqref{f:one-inflation} along a finite chain from $v_0$ to
$v$, each step satisfying $s<t\le2s$. There are
$O(1+\log(v/v_0))$ steps; covering larger balls by smaller balls at
each step gives the cell-count bound
\[
 \Lambda^{(k+1)v+O_k(v_0\log(v/v_0))+o(1)}
                (r_{l+v}/r_l)^{p_k}.
\]
Since $p_kx=(k+1)(1-d_0)$, this proves
\eqref{f:inflation-occupancy}. The finite mesh can include the whole
inflation chain from the outset, and $v_0/u$ can be as small as
needed before taking logarithmic limits.
\end{proof}

\begin{proposition}[Binary slope rigidity]\label{f:binary}
Every profile satisfies
\begin{equation}\label{f:binary-eq}
 f_i'(s)\in\{0,-k/2\}\quad\text{for almost every }s.
\end{equation}
\end{proposition}

\begin{proof}
Suppose a differentiability point, away from the transition seam,
has slope $-x$ with $0<x<k/2$. Choose a short rational block based at
$l$ very near this point and a finite relative mesh so that
$f(l+v)-f(l)=-xv+O(\eps u)$ throughout $[0,2u]$. Differentiability
allows errors arbitrarily small relative to every mesh spacing.
Put $d_0=1-2x/k\in(0,1)$ and apply the angular test with its fixed
positive exponent, together with the broad-witness test at a much
smaller $v_0$.

Condition on $\tau_l$ and $Q_{2u}$. Rescale this parent isotropically
to bounded size and write $\Delta=\Lambda^{-u}$. The information of
the fine isotropic cell is
\[
 \log_\Lambda(N_{2u}/N_u)+o(1).
\]
Projecting parallel to a drawn direction gives information
\[
 \log_\Lambda(N_{2u}/M_{l+u})+o(1).
\]
Its inverse-image cells, clipped to the parent, are comparable to
$T_{l+u}$: transverse width is $\Lambda^u$ and length is
$\Lambda^{2u}$. Thus Lemma~\ref{f:paths} gives the second formula.
Adding independent directions costs $o(1)$ information by Lemma
\ref{f:conditioning}, because all directions in $\tau_l$ drift by
at most the endpoint width on this box.

Draw $k+1$ directions conditionally independently. The angular test
shows, with probability tending to one, that their lifted vectors
form a basis with inverse norm at most $\Lambda^{O(\eps u)}$.
Indeed expose them successively; proximity to the span of the
previous ones is an affine hyperplane-strip event in the frequency
coordinates, and use one mesh precision much smaller than $u$.
In this basis each projection records all but one coordinate, up
to the same precision loss. Order the coordinates once and use
the chain rule for their information. The information of a
coordinate conditioned on all its predecessors is no larger,
up to the sampled $O(\eps u)$ slack of Lemma~\ref{f:conditioning},
than its information when one predecessor has been omitted.
Each coordinate is counted $k$ times in the $k+1$ projections.
This is the finite-partition Loomis--Whitney information inequality,
and gives
\[
 \log_\Lambda(N_{2u}/M_{l+u})
 \ge\frac{k}{k+1}\log_\Lambda(N_{2u}/N_u)-O(\eps u).
\]
The sample form follows from Markov's inequality for the finitely
many conditional probability ratios, as in the proof of Lemma
\ref{f:conditioning}; no uniform pointwise entropy inequality is
being asserted.

Let $g(v)=\log_\Lambda(N_v/M_l)$. The ratio definition gives
\[
 \log_\Lambda(M_{l+u}/M_l)
 =(k+2)u+q(f(l+u)-f(l))=(k+2)d_0u+O(\eps u).
\]
The preceding inequality becomes
\[
 g(2u)+kg(u)\ge(k+1)(k+2)d_0u-O(\eps u).
\]
Lemma~\ref{f:inflation} bounds $g(u)$ and $g(2u)$ above by
$(k+1)d_0u$ and $2(k+1)d_0u$, respectively. Their weighted sum is
exactly the last right side. Hence both upper bounds are equalities
up to $O(\eps u)$.

We now apply the robust discretized projection theorem in
Proposition~\ref{p:projection}, with ambient dimension $k+1$,
projection dimension $k$, and
$\alpha=(k+1)d_0\in(0,k+1)$. We give the set extraction because
exceptional points can depend on direction. Choose a cap and parent
box with conditional success probability close to one for all the
finite comparisons. Within this fixed cap and parent box, the
probabilities of typical $\Delta$ cells, without conditioning on the
drawn direction, lie between
$\Delta^{\alpha+O(\eps)}$ and $\Delta^{\alpha-O(\eps)}$.
Take their centers to form one set $\mathcal A$. It is
$\Delta$ separated and
\[
 \Delta^{-\alpha+O(\eps)}\le|\mathcal A|
                         \le\Delta^{-\alpha-O(\eps)}.
\]
For coarser balls, the mesh instances of
\eqref{f:inflation-occupancy}, divided by the saturated parent
occupancy, give the probability upper bounds with exponent
$\alpha$. Passing from probability to number by the preceding
individual-cell bounds gives precisely the nonconcentration
inequalities for $\mathcal A$, with $O(\eps)$ exponent tolerance.
Intermediate radii follow by nesting between consecutive mesh
values, whose ratio is chosen sufficiently fine on the logarithmic
scale.

Restrict the joint position-direction law first to successful
positions, then to directions for which conditional spatial and
projected success have probability bounded below. Starting with
success close to one makes both restrictions retain fixed positive
mass. Mixing the conditional angular laws preserves their
hyperplane nonconcentration, and these fixed-mass restrictions
change it by only a constant. For every retained direction, its
successful positions occupy at least
$\Delta^{O(\eps)}|\mathcal A|$ cells, because its conditional
individual-cell probabilities have the same upper bound. The
conditional projected information bound shows that these cells
project to at most
\[
 \Delta^{-kd_0-O(\eps)}
 =\Delta^{-k\alpha/(k+1)-O(\eps)}
\]
projection cells. The direction law obeys the required Schubert
nonconcentration: for lifted vectors $(2\xi,1)$, the relevant
Schubert test is exactly an affine hyperplane-strip test on $\xi$,
up to fixed constants. The theorem supplies a fixed positive
projection gain for every sufficiently large subset for some
retained directions, contradicting this upper bound once $\eps$
and the finite mesh tolerances are sufficiently small. All
nonconcentration exponents were fixed using $x>0$ before $\eps$
was sent to zero. This proves \eqref{f:binary-eq}.
\end{proof}

\subsection{The location of the full-slope interval}

Call a differentiability parameter a \emph{full-slope point} if
$f_i'=-k/2$, and a zero-slope point if $f_i'=0$.
The logarithmic derivative of $M_{i,s}$ is zero at the former and
$k+2$ at the latter.

\begin{lemma}[Midpoint closure]\label{f:midpoint}
For almost every pair $l<a$ of full-slope points in one field,
$(l+a)/2$ is a full-slope point. Consequently each field's
full-slope set is an interval modulo null sets. In the power case,
if $a$ is a full-slope point in one field, then $a/2$ is a
full-slope point in every other field, for almost every such $a$.
\end{lemma}

\begin{proof}
Use pairs for which all relevant parameters are differentiability
points, avoiding the one possible transition seam. Fubini's theorem
permits this for almost every pair. Set $b=(l+a)/2$ and choose
$0<\eta\ll\sigma\ll a-l$. By the angular test in a short block
above $l$, two independently drawn paths inside $\tau_l$ have
velocity difference $\Delta v$ satisfying
\[
 \kappa_i(l)^{-1}\Lambda^{-\eta}
 \lesssim|\Delta v|\lesssim\kappa_i(l)^{-1}
\]
with probability tending to one. Work inside the second path's
parent tube at $b+\sigma/2$. Adding the first direction costs no
logarithmic information by Lemma~\ref{f:conditioning}; the maximal
relative drift there is
\[
 \kappa_i(l)^{-1}\kappa_i(b+\sigma/2)^2
       =\kappa_i(a+\sigma),
\]
which is below the endpoint width when the block and $\sigma$ are
chosen in the interior. We used the exact scale identity
$\kappa_i(b)^2=\kappa_i(l)\kappa_i(a)$.

Suppose $b$ were a zero-slope point. The occupancy identity between
its scales $b-\sigma/2$ and $b+\sigma/2$ gives information
increment $(k+2)\sigma+o(\sigma)$. Spatial labels have at most
$k\sigma+o(\sigma)$ information, so the time label has at least
$2\sigma+o(\sigma)$.

On the other hand, a spatial intercept label for the first path
at scale $a-\sigma$ almost determines this time label. The spatial
uncertainty in the second parent is
$\kappa_i(b+\sigma/2)\ll\kappa_i(a-\sigma)$. Once the first
intercept is specified to width $\kappa_i(a-\sigma)$, the relative
motion determines time to accuracy at most
\[
 C\frac{\kappa_i(a-\sigma)}{|\Delta v|}
 \le C\kappa_i(l)\kappa_i(a-\sigma)\Lambda^\eta
 =C\kappa_i(b-\sigma/2)^2\Lambda^\eta.
\]
There are therefore at most $\Lambda^{\eta+o(1)}$ possible smaller
time labels. The whole second parent meets only boundedly many
first parents at $a+\sigma$, by the preceding drift identity and
the smaller time extent. Since $a$ is full slope, the first-path
information increment from $a+\sigma$ to $a-\sigma$ is
$o(\sigma)$. It remains an upper bound after conditioning on the
second parent, and the extra first-parent label has negligible
cost. The chain rule would bound the second time-label information
by $\eta+o(\sigma)$, contradicting $2\sigma+o(\sigma)$. First let
the sequence tend to its limit at fixed rational approximations,
then let $\eta/\sigma\to0$ and $\sigma\downarrow0$.

In the power case repeat the same geometry with the second path
in another field and with $l=0$. Root transversality gives
$|\Delta v|\ge\Lambda^{-o(1)}$, so the same time-recovery bound
holds and the drift is $O(h^{a+\sigma})$. The conditioning rule
applies to the other field's full weights. Full slope at $a$ in
the first field therefore forces full slope at $a/2$ in the
second field.

We finish with the measure-theoretic assertion. Let $E$ be a
measurable set satisfying the midpoint rule for almost every pair
of its points. For two density points $x,y$ choose equally short
intervals centered there on which $E$ has relative measure greater
than $9/10$. The convolution of their restricted indicators is
strictly positive on a neighborhood of $x+y$: after a small
translation the two high-density subsets still intersect in
positive measure. Fubini and the midpoint rule imply that a
neighborhood of $(x+y)/2$ belongs to $E$ modulo a null set. Taking
$x=y$ shows that the density points form an open set modulo null
sets; taking arbitrary $x,y$ shows its midpoint convexity. Repeated
midpoints, followed by openness, fill the interval between any
two such points. Thus $E$ is an interval modulo null sets.
\end{proof}

In the power case let $s_H=\lim\log_h H_i$, the same for every field
by \eqref{f:hbalance}. It is positive, since $A_i\ge h^c$.
The endpoint values of $f_i$ and the two allowed slopes show that
its full-slope interval has length $s_H$. If the left endpoints
of these intervals had positive minimum, the cross-field halving
rule would give a smaller left endpoint, a contradiction. An
interval beginning at zero then forces every other interval to
begin at zero by that rule. Therefore
\begin{equation}\label{f:power-end}
 f_i(s_H)=0\quad\text{for every }i.
\end{equation}

The analogous subpower conclusion needs control at the unbounded
left endpoint. We give the incidence calculation explicitly.

\begin{lemma}[Cross-field one-step inflation]
\label{f:cross-field-inflation}
Let $K\le R\le K^2$ be two inverse widths in a fixed finite packet
tree. At scale $K$, let $W_j$, $1\le j\le m=k+1$, be the original
selected packet fields testing the root high set, with no outer
smaller-inverse-width mask. Each $W_j$ is assembled from exactly its
selected scale-$R$ descendants through the frozen bounded masks and
scale-$K$ localizers; no omitted branch contributes to $W_j$.
The witnesses below test these same masked fields $W_j$.
Suppose their frequency labels,
including fixed enlargements, are $\delta$-transverse. Write
$\mathcal E_{j,K}=\sum_Ta_{j,T}^2w_T$ for their local square
envelopes and $\mathcal E_{j,R}$ for the envelopes of their
descendants at scale $R$ in the same frozen tree. Let $Y$ be a set
of distinct unit cubes with separate field witnesses at thresholds
$A_j$. Suppose that, on every $K$-ball meeting $Y$,
\[
 K^{k/2}\mathcal E_{j,K}(z)^{1/2}/A_j\le U,
 \qquad 1\le j\le m,
\]
at one, hence every, point $z$ in that ball, where $U\ge1$.
Then for every fixed $\eps>0$ and every $R$-ball centered at a
surviving root point $z$,
\begin{equation}\label{f:cross-field-inflation-bound}
 \#(Y\cap B_R(z))
 \lesssim_{\eps,\mathcal G}\delta^{-C_\eps}U^\eps R^{k+1}
 \prod_{j=1}^m\bigl(\mathcal E_{j,R}(z)/A_j^2\bigr)^{1/k}.
\end{equation}
The constant may depend on the fixed finite grid $\mathcal G$, but
not on the aspect ratio $R/K$.
\end{lemma}

\begin{proof}
Cover $B_R(z)$ by boundedly overlapping $K$-balls $D$. On one
meeting $Y$, the separate-unit-supremum form of
\Cref{p:local-multilinear}, with $p=p_k+\eps$, and comparability of
packet envelopes on $D$ give
\[
 \#(Y\cap D)\lesssim_\eps\delta^{-C_\eps}K^{kp/2}
   \prod_j\bigl(\mathcal E_{j,K}(z_D)/A_j^2\bigr)^{p/(2m)}.
\]
Since $kp_k/2=k+1$ and $p_k/(2m)=1/k$, this is
\[
 C_\eps\delta^{-C_\eps}K^{k+1}
 \prod_j\bigl(\mathcal E_{j,K}(z_D)/A_j^2\bigr)^{1/k}
 \left[K^{k/2}\prod_j
       \bigl(\mathcal E_{j,K}(z_D)/A_j^2\bigr)^{1/(2m)}
 \right]^\eps.
\]
The bracket is at most $U$, up to a fixed factor. Summing over
$D$ and replacing the sum by an integral uses only the fixed
variation of each envelope on a $K$-ball.

For these globally transverse groups, the local weighted cylinder
consequence of \Cref{p:kakeya} is
\[
 \frac1{|B_R|}\int_{B_R}\prod_j\mathcal E_{j,K}^{1/k}
 \lesssim\delta^{-C}
 \prod_j\left(\frac1{|CB_R|}
                 \int_{CB_R}\mathcal E_{j,K}\right)^{1/k}.
\]
To verify its normalization, expand the polynomial envelopes on
$B_R$ into summably weighted translated columns of radius $K$.
Their time length is at least $R$ because $R\le K^2$. A column
meeting the ball contributes a fraction comparable to $(K/R)^k$
of its coefficient to the average on a fixed enlargement. The
radius-$K$ multilinear Kakeya estimate contributes $K^{k+1}$;
the $m=k+1$ averages contribute $(R/K)^{k+1}$, giving exactly
$R^{k+1}$. The fixed reserve in the packet envelope power sums
the translated-column tails. This is the profile-free geometric
calculation in \eqref{f:one-inflation}, with the determinant factor
now $\delta^{-C}$.

For each field separately, the weighted spatial $L^2$ calculation
following \eqref{f:one-inflation} gives
\[
 \frac1{|CB_R(z)|}\int_{CB_R(z)}\mathcal E_{j,K}
 \lesssim_{\mathcal G}\mathcal E_{j,R}(z).
\]
Indeed reproduce the coarse weighted suprema, remove each
intervening bounded mask before separating its immediate child
caps by spatial almost orthogonality, and then use the final
$R$-packet envelope. In the final cap frame the normalized shift
has sizes $((K/R)(U_0+O(1)w),(K/R)^2w)$, so convolution preserves
the same fixed polynomial weight, as in the calculation preceding
\eqref{f:one-inflation} and following it. No power of $R/K$ is
lost. Combining these estimates proves
\eqref{f:cross-field-inflation-bound}.
\end{proof}

\begin{lemma}[Vanishing left limit]\label{f:left-limit}
In the subpower case,
\begin{equation}\label{f:left-limit-eq}
 \lim_{s\to-\infty}d_i(s)=0
 \quad\text{for every }i.
\end{equation}
Consequently $f_i(0)=0$, including when $e_i=0$.
\end{lemma}

\begin{proof}
Use the path law without an extra angular test, retaining each
path's attachment to an original surviving $H_i$ mask. Fix a field,
write $H=H_i$, $K=HS^s$, and let $\mu_z$ be the original local
direction law at a typical root point $z$. The clipped cap bound
and path domination give the slot estimate
\begin{equation}\label{f:slot}
 \mu_z(B(\omega,r))\lesssim S^{C+o(1)}r^k,
 \qquad r\ge K^{-1}.
\end{equation}
Indeed a slot at inverse width $K$ contains at most its frequency
volume times $H^k$, up to the appropriate bounded grid enlargement,
transition bins, each with clipped weight at most $H^{-k}S^{C_3}$;
path normalization loses only a subpower. The same bound summed
over slots proves the displayed formula.

There is also a lower bound at a typical sampled direction,
simultaneously at all relevant radii:
\begin{equation}\label{f:angular-typical-lower}
 \mu_z(B(\omega,r))\gtrsim S^{-1}r^k.
\end{equation}
To prove this without a number-of-radii loss, remove maximal
dyadic cubes on which the measure is less than $cS^{-1}$ times
volume. The maximal cubes are disjoint and lie in a bounded
frequency region, so their total measure is $O(S^{-1})$.
Use finitely many shifted dyadic grids to put a comparable cube
containing $\omega$ inside any ball in question. Every violating
ball is thereby charged to one of the removed cubes.

Let $y\in Y$ and $\tau=|t_y-t_z|$. The probability, under the
typical-direction restriction, that $y$ belongs to the sampled
enlarged $K$ tube is at most
\begin{equation}\label{f:smalltube-hit}
 S^{C+o(1)}(K/\tau)^k.
\end{equation}
Only $\tau\lesssim K^2S^{o(1)}$ matters. Hence the angular window
$K/\tau$ is no smaller than the slot precision up to a subpower,
and \eqref{f:slot} proves the bound.

If this event is nonempty, choose a typical direction $\omega_0$
realizing it. Every original path with direction within
$c\min(1,H/\tau)$ of $\omega_0$ has its attached enlarged $H$ mask
containing $y$. Those masks are already near $z$, the extra
transverse displacement is $O(H)$, and their enlarged duration
contains the relevant time interval. Thus
\eqref{f:angular-typical-lower} gives
\[
 \Pr_{\rm original}\{y\text{ hits the attached enlarged }H\text{ mask}\}
 \gtrsim S^{-1-o(1)}\min(1,H/\tau)^k.
\]
For $\tau\ge KS^{C_4}$, divide \eqref{f:smalltube-hit} by this
last bound. For $\tau\le H$ use $K/\tau\le S^{-C_4}$; for
$\tau\ge H$ use $K/H=S^s$. This gives the ratio bound
\[
 S^{C+1+o(1)}\big(S^{-kC_4}+S^{ks}\big).
\]
Sum this comparison \emph{once} over $y$. Every attached original
mask has occupancy at most $S^{C_1}$ by the deletion preceding
\eqref{f:ballcount}. Therefore
\[
 \mathbb E[\text{occupancy at time distance }\ge KS^{C_4}]
 \lesssim S^{C+C_1+1+o(1)}
                    (S^{-kC_4}+S^{ks})=o(1)
\]
for sufficiently large fixed $C_4$ and $s<-2C_4$. There is no
iteration over a number of intermediate scales.

The occupancy identity \eqref{f:occupancy-law}, together with the
typical leaf probability comparisons, turns its effective
occupancy into an ordinary count up to $S^{o(1)}$. With probability
tending to one there are no far points, so a ball of radius
$O(KS^{C_4})$ contains at least
\begin{equation}\label{f:left-lower}
 M_{i,s}S^{-o(1)}=S^{qd_i(s)-o(1)}
\end{equation}
root points. The identity is exact on the logarithmic scale:
\[
 M_{i,s}=(H_iS^s)^{k+2}
             (S^{f_i(s)}/H_i^{k/2})^q=S^{qd_i(s)},
\]
because $kq/2=k+2$.

For the upper bound use a common ball scale at
$s'=s+C_4+1<0$. The inverse widths in the fields differ by only
subpowers, by \eqref{f:subbalance}. Choose a common first width below
all field-specific first widths and a common terminal width above all
field-specific terminal widths, each within an $S^{o(1)}$ factor.
Inserting these levels into the fixed grid preserves the tested field;
the local polynomial-weight convolution and bounded-multiplicity
comparison in \Cref{p:hierarchy} change its square-envelope bounds by
at most $S^{o(1)}$. Begin at
the first gap-edge parameter $s_{\min}\le s$ of a fixed finite
grid, where no earlier smaller-inverse-width masks intervene.
Write $\mathcal E_{j,s}(z)=\sum_Ta_{j,T}^2w_T(z)$ for the local
square envelope of the original selected outputs at that layer;
this is distinct from the integrated energy $E_{j,s}$ in
\eqref{f:Esat}. On every first-edge ball carrying witnesses, the
selected ratio class and clipped square give
$\kappa(s_{\min})^{k/2}\mathcal E_{j,s_{\min}}^{1/2}/A_j\le S^C$:
the remaining exponent is $d_j(s_{\min})\le ke_j/2$.
The transversality loss is $S^{o(1)}$. We may apply
\Cref{f:cross-field-inflation} once from this first edge to $s'$,
because for every fixed finite range
\[
 \frac{\kappa(s_{\min})^2}{\kappa(s')}
 =H_iS^{2s_{\min}-s'}\longrightarrow\infty.
\]
Here $H_i$ grows faster than every fixed power of $S$. The single
step keeps the determinant power and the $\eps$ cost independent
of grid length. At the root-centered terminal ball,
$\mathcal E_{j,s'}\le S^{o(1)}r_{j,s'}^2$ by the clipped square
bound. Substitution in \eqref{f:cross-field-inflation-bound} gives the count
upper bound
\begin{align}
 S^{o(1)+C\eps}\kappa(s')^{k+1}
           \prod_j(\mathcal E_{j,s'}/A_j^2)^{1/k}
 &\le S^{(k+1)s'+(2/k)\sum_jf_j(s')+o(1)+C\eps}\notag\\
 &=S^{\frac{p_k}{k+1}\sum_jd_j(s')+o(1)+C\eps}.
                                                        \label{f:left-upper}
\end{align}
The $H_j$ factors cancel since their ratios are subpowers.
Take limits for this finite comparison and then send $\eps$ to
zero. Equations~\eqref{f:left-lower} and \eqref{f:left-upper} yield
\[
 qd_i(s)\le\frac{p_k}{k+1}\sum_jd_j(s+C_4+1).
\]

By \eqref{f:d-bounded}, every $d_i$ is bounded, nonnegative, and
nondecreasing. Let $L_i=\lim_{s\to-\infty}d_i(s)$. Sending $s$ to
$-\infty$ and summing the preceding inequalities gives
\[
 (q-p_k)\sum_iL_i\le0.
\]
Since $q>p_k$ and $L_i\ge0$, all the limits vanish.

The full-slope interval must extend to the unbounded left endpoint:
otherwise $f_i$ is constant sufficiently far to the left and
$d_i(s)=f_i(s)+ks/2$ becomes negative. On this interval $d_i'=0$,
so the just proved limit gives $d_i=0$. After its right endpoint
$f_i$ is constant, and its fine endpoint value makes that constant
zero. Continuity at the right endpoint therefore places it at
$s=0$. If the full-slope interval fills the domain, the endpoint
value instead gives $e_i=0$. In either case $f_i(0)=0$.
\end{proof}

The profiles forced by these arguments are shown in
\Cref{f:profiles-figure}.
\begin{figure}[htbp]
 \centering
 \begin{tikzpicture}[x=1cm,y=1cm,
  axis/.style={-{Stealth[length=1.6mm]},black!65,line width=.5pt},
  profile/.style={blue!65!black,line width=1.2pt},
  every node/.style={font=\small}]
 \begin{scope}
  \node at (1.65,3.15) {Power case: $\Lambda=h$};
  \node at (1.65,2.72) {$\kappa_i(s)=h^s$};
  \draw[axis] (-.15,0)--(3.85,0) node[right] {$s$};
  \draw[axis] (0,-.15)--(0,2.45) node[above left] {$f_i(s)$};
  \draw[profile] (0,2.1)--(2.25,0)--(3.5,0);
  \fill[blue!65!black] (0,2.1) circle (1.6pt)
                         (2.25,0) circle (1.6pt)
                         (3.5,0) circle (1.6pt);
  \node[below] at (0,-.06) {$0$};
  \node[below] at (2.25,-.06) {$s_H$};
  \node[below] at (3.5,-.06) {$1$};
  \node[above right] at (.85,1.38) {slope $-k/2$};
  \node at (1.65,-.75) {$f_i(s)=\dfrac{k}{2}(s_H-s)_+$};
  \node at (1.65,-1.35) {$0\le s\le1,\quad0<s_H\le1$};
 \end{scope}
 \begin{scope}[xshift=8.15cm]
  \node at (0,3.15) {Subpower case: $\Lambda=S$};
  \node at (0,2.72) {$\kappa_i(s)=H_iS^s$};
  \draw[axis] (-2.6,0)--(2.05,0) node[right] {$s$};
  \draw[axis] (0,-.15)--(0,2.45);
  \node[anchor=west] at (.2,2.1) {$f_i(s)$};
  \draw[profile,{Stealth[length=2mm]}-] (-2.25,2.1)--(0,0)--(1.5,0);
  \fill[blue!65!black] (0,0) circle (1.6pt)
                         (1.5,0) circle (1.6pt);
  \node[below] at (0,-.06) {$0$};
  \node[below] at (1.5,-.06) {$e_i$};
  \node at (-1.3,2.3) {slope $-k/2$};
  \node at (0,-.75) {$f_i(s)=(-ks/2)_+$};
  \node at (0,-1.35) {$-\infty<s\le e_i,\quad e_i\ge0$};
 \end{scope}
\end{tikzpicture}

 \caption{The forced amplitude-ratio profiles
 $r_{i,s}=\Lambda^{f_i(s)+o(1)}$. Their shapes follow from
 \Cref{f:profiles,f:binary,f:midpoint}, \eqref{f:power-end}, and
 \Cref{f:left-limit}. The panels show generic cases; the terminal
 flat segment may collapse at $s_H=1$ or $e_i=0$, respectively.
 The subpower domain continues to the left without a finite endpoint.
 The recomputed transition ratio is $\Lambda^{o(1)}$ in both cases,
 as used in the final incidence contradiction.}
 \label{f:profiles-figure}
\end{figure}

\subsection{The final incidence contradiction}

We finish the proof of Theorem~\ref{f:frame}. In the power case
\eqref{f:power-end} says that the recomputed transition ratio is
$\Lambda^{o(1)}$; in the subpower case the same assertion follows
from Lemma~\ref{f:left-limit}. The original transition masks have
remained fixed throughout.

One final occupancy deletion makes the use of packet tails
independent of any convergence rate for the profiles. Fix an
arbitrarily small $\eta>0$ and enlarge transition cells by
$\Lambda^\eta$ in their native coordinates. In the power case
count all visits. In the subpower case count points with
$e_b^2(z)\le CM_c\Lambda^{C_0\eta}$, where $C_0$ is sufficiently
large for polynomial weight comparison. This includes every
neighbor of a surviving original mask. The incident squares sum
to at most $\Lambda^{C_2\eta+o(1)}|Y|$. Since the transition ratio
is $\Lambda^{o(1)}$ and
$H_i^{k+2}A_i^{-q}=1$, the weighted $q$ charge of cells with more
than $\Lambda^{C_1'\eta}$ visits is at most
\[
 \Lambda^{(C_2-C_1')\eta+o(1)}D_i.
\]
Choose $C_1'>C_2$ sufficiently large. Lemma~\ref{f:selection}
propagates this error to the root with a negative power saving.
Do not reselect ratios after this deletion: all other transition
fields remain the same. The fixed packet decay, with the square
bound at surviving masks, makes tails outside the chosen radius
negligible relative to the high threshold. Thus each remaining
root point receives a substantial contribution from nearby
transition masks, each visited at most
$\Lambda^{C_1'\eta}$ times.

The field threshold is $A_i\Lambda^{-o(1)}=H_i^{k/2}\Lambda^{-o(1)}$.
At a root point the square sum of the transition backgrounds is
$\Lambda^{o(1)}$, and each transition amplitude is at most its
background times $\Lambda^{o(1)}$. Weighted Cauchy and bounded
absolute sums of the ancestor cutoffs therefore require at least
$H_i^k\Lambda^{-o(1)}$ nearby active masks. Letting the fixed
$\eta$ be arbitrarily small, double counting visits gives at least
\[
 |Y|H_i^k\Lambda^{-o(1)}
\]
distinct original transition masks.

Choose a field with $S_i=S$. Each of these masks had original
amplitude $a_T^{\mathrm{orig}}\ge cSH_i^{-k/2}$ in
\eqref{f:original-amplitudes}, and their original integrated square
budget is bounded by \eqref{f:original-energy}. Hence
\[
 D_i\gtrsim H_i^{k+2}
       \big(|Y|H_i^k\Lambda^{-o(1)}\big)
       \big(S^2H_i^{-k}\big)
 =|Y|H_i^{k+2}S^2\Lambda^{-o(1)}.
\]
This contradicts \eqref{f:hbalance} in the power case, because
$S$ is a fixed positive power of $\Lambda=h$, and contradicts
\eqref{f:subbalance} in the subpower case, because $S=\Lambda$.
The exact-packet estimate follows. Lemma~\ref{f:exact} transfers it
to compact profiles and fixes a finite derivative order after all
the fixed weight margins have been chosen. The countersequence
construction proves the uniform quantitative form
$C\delta^{-C}$ in \eqref{f:mainbound}, completing the proof.
